\chapter{Implementación}
\label{chap:implementacion}

\ph{Este capítulo describe \textbf{cómo} se ha materializado el diseño en código:
qué tecnologías se han usado, qué módulos se han implementado y qué decisiones
concretas se han tomado en cada uno.}

\section{Análisis de tecnologías}
\label{sec:analisis_tecnologias}

\ph{Justifica la elección de las tecnologías principales del proyecto.
Si has comparado varias opciones, hazlo explícito (es más convincente que
afirmar la elegida sin alternativas).}

\subsection*{\ph{Necesidad técnica 1}}

\ph{Plantea el problema concreto que motivó la elección
(p.\,ej., procesar archivos en formato X, validar reglas, automatizar Y\dots)
y qué características debe tener la tecnología elegida.}

\subsection*{\ph{Necesidad técnica 2}}

\ph{Otra necesidad que orienta una decisión tecnológica.}

\subsection*{\ph{Necesidad técnica 3}}

\ph{Otra necesidad que orienta una decisión tecnológica
(p.\,ej., experiencia de usuario, despliegue, rendimiento, mantenibilidad).}

\subsubsection*{Comparativa de tecnologías del backend}

\begin{itemize}
  \item \textbf{\ph{Tecnología A}}: \ph{descripción y ventajas/inconvenientes.}
  \item \textbf{\ph{Tecnología B}}: \ph{descripción y ventajas/inconvenientes.}
  \item \textbf{\ph{Tecnología C}}: \ph{descripción y ventajas/inconvenientes.}
\end{itemize}

\subsubsection*{Comparativa de tecnologías del frontend}

\begin{itemize}
  \item \textbf{\ph{Tecnología A}}: \ph{descripción y ventajas/inconvenientes.}
  \item \textbf{\ph{Tecnología B}}: \ph{descripción y ventajas/inconvenientes.}
  \item \textbf{\ph{Tecnología C}}: \ph{descripción y ventajas/inconvenientes.}
\end{itemize}

\subsubsection*{Análisis de la base de datos}

\begin{enumerate}
  \item \textbf{\ph{Opción A}}: \ph{justificación.}
  \item \textbf{\ph{Opción B}}: \ph{justificación.}
\end{enumerate}

\subsubsection*{Plataformas de despliegue}

\begin{itemize}
  \item \textbf{\ph{Plataforma A}}: \ph{descripción y casuística.}
  \item \textbf{\ph{Plataforma B}}: \ph{descripción y casuística.}
  \item \textbf{\ph{Plataforma C}}: \ph{descripción y casuística.}
\end{itemize}

\subsubsection*{Tecnologías finales seleccionadas}

\begin{itemize}
  \item \textbf{Desarrollo}: \ph{tecnologías elegidas para backend / frontend.}
  \item \textbf{Base de datos}: \ph{motor elegido.}
  \item \textbf{Contenedor}: \ph{tecnología de contenerización (Docker, Podman\dots).}
  \item \textbf{Despliegue}: \ph{plataforma final.}
\end{itemize}

\section{Tecnologías utilizadas}
\label{sec:tecnologias}

\ph{Tras la justificación, describe brevemente cómo se usa cada tecnología
en el proyecto. Una tecnología por párrafo, indicando módulo o archivo
en que aparece.}

\ph{Sugerencia de bloques:}

\begin{itemize}
  \item \ph{Backend (framework, gestión de dependencias, configuración).}
  \item \ph{Autenticación (estrategia, almacenamiento de credenciales,
        gestión de sesión o token).}
  \item \ph{Persistencia (ORM, motor, migraciones).}
  \item \ph{Lógica de dominio (módulos en los que se divide).}
  \item \ph{Generación de artefactos / motor de plantillas (si aplica).}
  \item \ph{Capa de presentación (HTML/CSS/JS o framework).}
  \item \ph{Infraestructura de desarrollo (Docker, Compose, Make\dots).}
  \item \ph{Suite de pruebas (framework, marcadores, comandos).}
  \item \ph{Automatización CI/CD (proveedor y workflows principales).}
\end{itemize}

\section{\ph{Módulo de validación / lógica principal}}
\label{sec:modulo_principal}

\ph{Describe el módulo central de la lógica de negocio. Suele ser el que aplica
las reglas del dominio (validador, motor de reglas, planificador\dots).}

\ph{Cubre:}
\begin{itemize}
  \item \ph{Función o clase principal y su contrato (entrada/salida).}
  \item \ph{Estructura del resultado (campos, tipos, semántica).}
  \item \ph{Cómo lo consumen otros módulos (API, generadores\dots).}
\end{itemize}

\subsection{\ph{Validación sintáctica / comprobaciones básicas}}
\label{subsec:validacion_sintactica}

\ph{Comprobaciones de presencia de campos, pertenencia a conjuntos enumerados,
tipos de datos\dots En definitiva, lo que se puede verificar sin entrar en
la semántica del dominio.}

\subsection{\ph{Validación semántica / restricciones de dominio}}
\label{subsec:validacion_semantica}

\ph{Reglas de dominio que combinan varios campos (rangos, dependencias,
fórmulas, exclusiones). Documenta cada regla con su descripción y, si aplica,
su expresión matemática.}

% Hueco para la expresión matemática de una regla. Descomenta y adapta:
% \[
% \text{regla}_1: \quad f(x_1, x_2) \geq c_1
% \]

\ph{(Espacio reservado para la expresión matemática de las reglas que la
necesiten, p.~ej.\ $\mathrm{regla}_1: f(x_1, x_2) \geq c_1$. Descomenta y
adapta el bloque comentado del fuente.)}

\section{\ph{Módulo de generación de artefactos}}
\label{sec:generador}

\ph{Describe el componente que produce el artefacto final
(documento JSON, código fuente, fichero CAD, informe PDF\dots) a partir de
la representación intermedia.}

\subsection{\ph{Cálculos o transformaciones específicas}}
\label{subsec:calculos}

\ph{Si la generación implica cálculos no triviales (geometría, fórmulas
analíticas, transformaciones de coordenadas\dots), descríbelos con sus
fórmulas o pseudocódigo.}

\subsection{\ph{Síntesis del artefacto final}}
\label{subsec:sintesis}

\ph{Describe el proceso de generación final paso a paso, separando
preparación de datos, renderizado de plantillas y operaciones específicas
de la herramienta destino.}

\subsection{\ph{Operaciones complementarias}}
\label{subsec:operaciones}

\ph{Si tras la generación hay operaciones complementarias
(ensamblado, postprocesado, decoración con metadatos\dots), descríbelas aquí.}

\section{API REST / capa de servicios}
\label{sec:api}

\ph{Describe la API expuesta por el sistema: agrupa los \emph{endpoints} por
bloque temático y, para cada endpoint, indica método, ruta, autenticación
requerida y semántica básica.}

\subsection{Endpoints implementados}
\label{subsec:endpoints}

\ph{Bloques sugeridos:}

\ph{Páginas / vistas:}
\begin{itemize}
  \item \ph{\texttt{GET /ruta\_publica}}: \ph{descripción.}
  \item \ph{\texttt{GET /ruta\_protegida}}: \ph{descripción.}
\end{itemize}

\ph{Endpoints de lógica principal:}
\begin{itemize}
  \item \ph{\texttt{POST /api/\dots}}: \ph{descripción y respuestas (200, 4xx, 5xx).}
  \item \ph{\texttt{GET  /api/\dots}}: \ph{descripción.}
\end{itemize}

\ph{Endpoints de autenticación:}
\begin{itemize}
  \item \ph{\texttt{POST /api/auth/register}}: \ph{descripción.}
  \item \ph{\texttt{POST /api/auth/login}}: \ph{descripción.}
  \item \ph{\texttt{POST /api/auth/logout}}: \ph{descripción.}
\end{itemize}

\ph{Endpoints de gestión de recursos del usuario:}
\begin{itemize}
  \item \ph{\texttt{POST /api/\dots}}: \ph{descripción.}
  \item \ph{\texttt{GET  /api/\dots}}: \ph{descripción.}
\end{itemize}

\subsection{Persistencia y recuperación}
\label{subsec:persistencia}

\ph{Explica cómo se guardan y recuperan los recursos principales: qué se
persiste, qué validaciones se aplican antes de persistir y cómo se garantiza
el aislamiento entre usuarios.}

\subsection{\ph{Descarga / exportación de artefactos}}
\label{subsec:descarga}

\ph{Si el sistema permite descargar artefactos generados, describe los
escenarios (descarga inmediata tras generación, descarga desde el historial
de configuraciones guardadas\dots) y la estrategia de regeneración bajo demanda
frente a almacenamiento permanente.}

\section{Integración continua}
\label{sec:ci}

\ph{Describe la automatización de calidad y despliegue del proyecto.
Si usas un proveedor concreto (GitHub Actions, GitLab CI, Jenkins\dots),
nómbralo y comenta los workflows configurados.}

\ph{Workflows típicos:}

\begin{itemize}
  \item \textbf{\ph{Validación de commits}}: \ph{convención de mensajes, autores\dots}
  \item \textbf{\ph{Linting / análisis estático}}: \ph{herramienta y comandos.}
  \item \textbf{\ph{Pruebas}}: \ph{ejecución de la suite en cada \textit{push} y \textit{PR}.}
  \item \textbf{\ph{Cobertura}}: \ph{umbral mínimo y políticas de fallo.}
  \item \textbf{\ph{Release}}: \ph{condiciones para crear releases.}
  \item \textbf{\ph{Despliegue}}: \ph{construcción y publicación de imagen / artefacto.}
\end{itemize}

\ph{Si dispones de comandos locales equivalentes a los workflows remotos
(Make, scripts\dots), enuméralos para que el lector pueda reproducir
las comprobaciones.}

% --- Transición al siguiente capítulo ---
\ph{Cierra el capítulo con un párrafo puente de 2--4 líneas: qué queda
establecido aquí y qué pregunta abre el capítulo siguiente. Por ejemplo:
«Con el sistema implementado y la integración continua en marcha, falta
demostrar que se comporta como exigen los requisitos del
Capítulo~\ref{chap:requisitos}. El Capítulo~\ref{chap:pruebas} presenta la
estrategia de pruebas con la que se verifica».}
